%% tags: [loop invariants]
%% source: 2023-fa-redemption_midterm_01
\begin{prob}
    Consider this code which partitions a list in a special way:

    \inputminted[lastline=15]{python}{code.py}

    Which of the following loop invariants are true for this code? Select all
    that apply.

    Note: any statement about an empty set or list is considered to be automatically true.

    \begin{choices}[rectangle]
        \correctchoice After each iteration of the \mintinline{python}{for} loop, all numbers in
            \mintinline{python}{numbers[:barrier_ix]} are even.
        \choice After each iteration of the \mintinline{python}{for} loop, all numbers in
            \mintinline{python}{numbers[barrier_ix:]} are odd.
        \correctchoice After each iteration of the \mintinline{python}{for} loop,
            \mintinline{python}{numbers[0]} is greater than or equal to all numbers in
            \mintinline{python}{numbers[:barrier_ix]}.
        \choice After each iteration of the \mintinline{python}{for} loop,
            \mintinline{python}{numbers[0]} is the maximum of \mintinline{python}{numbers[:i+1]}.
    \end{choices}

    \begin{soln}
        1st option: After each iteration of the \mintinline{python}{for} loop, all numbers in
            \mintinline{python}{numbers[:barrier_ix]} are even.

        3rd option: After each iteration of the \mintinline{python}{for} loop,
        \mintinline{python}{numbers[0]} is greater than or equal to all numbers in
        \mintinline{python}{numbers[:barrier_ix]}.
    \end{soln}

\end{prob}
